\section{The complete marked metric and its attained relative quotient}


\subsection{Source, coefficients, and exact scope}

The human source is Alain Connes and Caterina Consani,
\emph{Knots, primes and class field theory},
\href{https://arxiv.org/abs/2501.06560v1}{arXiv:2501.06560v1}.
The original author file \path{arxiv-pisa.tex}, lines 880--1035,
was read, including Section 4 and the original labels
\texttt{localv}, \texttt{structure}, and \texttt{structure1}.
These supply the actual semilocal transition maps, closed stalks,
generic Bruhat--Schwartz coefficients, and rational crossed-product
labels used below. No replacement of the coefficient sheaf is made.
The independent calculation verifies (CCF8)--(CCF19) and
(CCF26)--(CCF34) in the local programme source
\path{CONNES_SUPPORT_FUNCTOR_AND_COHOMOLOGY.tex}.
The formulas below prove the cohomology, closure, and positive
receiver directly. Their assertions concern underlying complex
vector sheaves and the explicitly defined Hilbert spaces.

Put $X=\Spec\Z$, $V=\mathcal S(\mathbb A_\Q)$, and
$B=\bigoplus_{r\in\Q^\times}VU_r$, retaining its original
crossed-product multiplication when an algebra is required.
For each finite prime $p$, let $\xi_p=1_{\Z_p}$ and let
$V^{(p)}$ contain all factors other than $p$, including the real
factor. The original stalks are
\[
 V_p=\C\xi_p\otimes V^{(p)},\qquad
 B_p=\bigoplus_{r:\,v_p(r)=0}V_pU_r.
 \tag{MQ1}\label{mq:eq:stalks}
\]
They follow directly by taking the union of all semilocal stages
whose excluded finite-prime set does not contain $p$. The other
finite factors and the rational label are unrestricted away from
$p$, while the displayed conditions persist at $p$.

There are two versions of the calculation. In the coefficient
version let $W=V$ and $W_p=V_p$. In the labelled version let
$W=B$ and $W_p=B_p$. Write $\mathcal H$ for the corresponding
original sheaf and set
\[
 D_W=\bigoplus_p W/W_p,\qquad
 \delta w=([w]_p)_p.
 \tag{MQ2}\label{mq:eq:marked}
\]
The sum is finite for every $w$: each adelic test has only finitely
many nonstandard finite factors and each rational label has only
finitely many nonzero valuations. In particular this construction
retains every prime mark rather than identifying equal coefficients
at different marks.

\subsection{The sheaf cohomology and its full rational labels}

\begin{theorem}[The actual arithmetic resolution]
\label{mq:thm:resolution}
There is an exact sequence of sheaves
\[
 0\longrightarrow\mathcal H\longrightarrow c_XW
 \xrightarrow{\delta}\bigoplus_p i_{p*}(W/W_p)
 \longrightarrow0,
 \tag{MQ3}\label{mq:eq:resolution}
\]
whose two right-hand sheaves are flasque. Consequently
\[
 R\Gamma(X,\mathcal H)=[W\xrightarrow{\delta}D_W],
 \qquad H^1(X,\mathcal H)=D_W/\delta W,
 \qquad H^n(X,\mathcal H)=0\quad(n\ge2).
 \tag{MQ4}\label{mq:eq:cohomology}
\]
In degree zero the groups are respectively $\mathcal S(\mathbb R)$ and
$\mathcal S(\mathbb R)\rtimes\{\pm1\}$, as in the human source.
\end{theorem}
\begin{proof}
On $U=X\setminus F$, membership in the kernel means that the
coefficient has the factor $\xi_p$ at every $p\notin F$ and,
in the labelled case, that each surviving label is a unit at those
primes. The source's local factorization lemma proves the first
condition: support in $\Z_p$ together with invariance under its
additive translations is equivalent to
$f(x_p,y)=\xi_p(x_p)f(0,y)$. Applying this to every retained prime
gives exactly the original semilocal sections. Independence of the
group-algebra labels proves the assertion coefficient by coefficient.
At a closed point the quotient map is onto by definition; at the
generic point the target stalk is zero. This proves sheaf exactness.

Sections of the direct sum on a nonempty open have finite support.
Indeed their generic germ is zero, so the section vanishes on a
smaller cofinite open. Its remaining support is finite. Conversely
any finite collection of marked values gives a section, by the
skyscraper descriptions. Restrictions delete excluded marks and
are onto. This proves flasqueness. Every nonempty open of $X$ is
irreducible, so the constant sheaf has value $W$ there and identity
restrictions, hence is flasque too. Applying global sections to
this flasque resolution proves \eqref{mq:eq:cohomology}. Its kernel
consists of tests standard at every finite prime; a label standard
at every prime is precisely $1$ or $-1$. This proves the last claim.
\end{proof}

Every local algebraic splitting
$\mathcal S(\Q_p)=\C\xi_p\oplus E_p$ expands the original restricted
tensor product as the algebraic direct sum
\[
 V=\bigoplus_{T\text{ finite}}V_T,\qquad
 V_T=\mathcal S(\mathbb R)\otimes\bigotimes_{p\in T}E_p,
 \tag{MQ5}\label{mq:eq:algebraic-channels}
\]
with standard factors outside $T$. This follows by expanding each
finite tensor. Uniqueness follows by applying the two complementary
local projections within any finite set containing the factors of
a proposed finite relation.

For uniform notation, in the unlabelled version a channel $a$ is
$T$ and its marked set is $J_a=T$. In the labelled version it is
$a=(r,T)$ and
\[
 D(r)=\{p:v_p(r)\ne0\},\qquad J_a=T\cup D(r),
 \qquad W_a=V_TU_r.
 \tag{MQ6}\label{mq:eq:marked-set}
\]
In the unlabelled version set $W_a=V_T$. A channel survives in
$W/W_p$ exactly when $p\in J_a$. Thus the entire arithmetic
complex, with no rational labels removed, is
\[
 [W\longrightarrow D_W]
   =\bigoplus_a[W_a\xrightarrow{\Delta}W_a^{J_a}].
 \tag{MQ7}\label{mq:eq:channel-complex}
\]
An empty marked set gives zero target. Taking this finite diagonal
quotient in every channel proves
\[
 H^1(X,\mathcal H)
   =\bigoplus_{a:\,|J_a|\ge2}
       W_a\otimes(\C^{J_a}/\Delta\C).
 \tag{MQ8}\label{mq:eq:full-mixed-cohomology}
\]
The evaluation complement $E_p=\ker(f\mapsto f(0))$ gives exactly
the original (CCF13)--(CCF15). Nothing in this computation requires
it to be preserved by Fourier transform or by nonunit scaling.

\subsection{The full support extension and the arithmetic image}

Let $S=\SpecLat L$ for the original finite nontrivial bounded
distributive lattice, and let $Y=S\triangleright X$. Its opens are
the opens of $S$ and $S\cup U$ for nonempty arithmetic opens $U$.
Use the universal sheaf $U_S\mathcal H$, whose arithmetic component
is $\mathcal H$, whose support component is $c_SW$, and whose
gluing map is the original generic coefficient, constant on $S$.
Write $C^\bullet(S,W)$ for the strict-chain finite cochains, with
the original lower-open specialization convention. Its differential
is alternating deletion of chain entries. The complete finite
cochain resolution and recollement are proved in
\path{finite_support_recollement.tex}, (FL8)--(FL22).

The actual relative complex is
\[
 \begin{gathered}
 K^0=W,\qquad K^1=D_W\oplus C^0(S,W),\qquad
 K^n=C^{n-1}(S,W)\quad(n\ge2),\\
 d^0w=(\delta w,\Delta w),\qquad
 d^1(d,c)=d_Sc,\qquad d^n=d_S\quad(n\ge2).
 \end{gathered}
 \tag{MQ9}\label{mq:eq:relative-complex}
\]
Indeed \eqref{mq:eq:resolution} gives the arithmetic complex and
the generic morphism is the identity on its degree-zero $W$ and
zero on its degree-one $D_W$. The mapping fibre of the map to
$C^\bullet(S,W)$ has precisely these terms and, after multiplying
successive support coordinates by alternating signs, these
differentials. This fixes both its coefficient map and its sign.

\begin{proposition}[The image in the original arithmetic group]
\label{mq:prop:arithmetic-image}
The full relative cohomology has the exact sequence
\[
 0\longrightarrow D_W\xrightarrow{d\mapsto[(d,0)]}
 H_X^1(Y,U_S\mathcal H)
 \longrightarrow W\otimes\widetilde H^0(S,\C)
 \longrightarrow0,
 \tag{MQ10}\label{mq:eq:relative-exact}
\]
and $H_X^{j+1}=W\otimes H^j(S,\C)$ for $j\ge1$.
The actual map from relative to arithmetic cohomology is
\[
 [(d,c)]\longmapsto[d]\in D_W/\delta W.
 \tag{MQ11}\label{mq:eq:arithmetic-map}
\]
Its restriction to the canonical $D_W$ is the quotient map and
is surjective. Hence (CCF19) has exactly the claimed arithmetic
image \eqref{mq:eq:full-mixed-cohomology}.
\end{proposition}
\begin{proof}
A closed degree-one representative has $d_Sc=0$, so $c$ is
constant on each support component. Modulo the global diagonal
its class gives the right-hand map of \eqref{mq:eq:relative-exact}.
It is onto. If $c=\Delta w$, subtract the boundary
$(\delta w,\Delta w)$ to obtain $(d-\delta w,0)$. This proves
exactness. The indicated injection has zero kernel because
$(d,0)=(\delta w,\Delta w)$ implies $w=0$. All higher groups
follow directly from the remaining support cochains. Projection
of \eqref{mq:eq:relative-complex} onto $[W\to D_W]$ is a chain map:
it is the identity on $W$, sends $(d,c)$ to $d$, and is zero
thereafter. This proves \eqref{mq:eq:arithmetic-map}, including its
well-definedness and surjectivity on the canonical summand.
\end{proof}

Choosing a point $s_0$ gives the algebraic retraction
$[(d,c)]\mapsto d-\delta c_{s_0}$. It is well defined because
adding $(\delta w,\Delta w)$ adds equal terms. It is the identity
on the canonical $D_W$. The choice is explicit; the injection
and the arithmetic quotient in \eqref{mq:eq:relative-exact}--
\eqref{mq:eq:arithmetic-map} are canonical.

\subsection{Orthogonal unramified projections and the attained quotient}

Use additive Haar measure with $\operatorname{vol}(\Z_p)=1$ and
the real measure $dx$. Set $H=L^2(\mathbb A_\Q)$ and
\[
 P_pf=\xi_p\otimes\int_{\Z_p}f(t,\cdot)\,dt,
 \qquad Q_p=I-P_p.
 \tag{MQ12}\label{mq:eq:orthogonal-projection}
\]
This is the orthogonal projection onto
$\C\xi_p\widehat\otimes H^{(p)}$, since $\|\xi_p\|=1$.
It preserves adelic tests and its image on them is exactly $V_p$.
The projections for different primes commute because they act on
different tensor factors. Use from now on
\[
 E_p^\perp=\{f\in\mathcal S(\Q_p):\int_{\Z_p}f=0\},
 \quad V_T^\perp=\mathcal S(\mathbb R)\otimes
                  \bigotimes_{p\in T}E_p^\perp.
 \tag{MQ13}\label{mq:eq:orthogonal-channels}
\]
The same proof as \eqref{mq:eq:algebraic-channels} gives an algebraic
direct sum of mutually orthogonal spaces. Let $H_T$ be the Hilbert
closure of $V_T^\perp$ in $H$. Their orthogonal sum is all of $H$:
finite adelic tensors are dense, and their finite orthogonal
expansions lie in these summands. This also proves that each local
orthogonal test space is dense in its local Hilbert complement,
by applying $I-P_p$ to a dense family of local tests.

In the labelled case use
$H_B=\widehat\bigoplus_r HU_r$, with counting measure on the
original rational labels. Put
\[
 Q_{p,r}=\begin{cases}I-P_p,&v_p(r)=0,\\I,&v_p(r)\ne0.
 \end{cases}
 \tag{MQ14}\label{mq:eq:labelled-projection}
\]
It realizes the quotient by $B_p$ label by label. On $H_TU_r$
it is the identity exactly for $p\in T\cup D(r)$ and zero
otherwise. In the unlabelled case use $Q_p$ and omit $r$.
Consequently the Hilbert completion of the marked quotient is
\[
 \overline D_W=
 \widehat\bigoplus_{a:\,k_a>0} H_a\otimes\C^{J_a},
 \qquad k_a=|J_a|,
 \tag{MQ15}\label{mq:eq:marked-hilbert}
\]
where $H_a=H_T$ or $H_TU_r$. All finite-factor and rational-label
marks are included in $J_a$ exactly as in \eqref{mq:eq:marked-set}.

\begin{theorem}[The closure and the injective positive arithmetic receiver]
\label{mq:thm:closure}
The closure of the actual generic diagonal is
\[
 \overline{\delta W}=
 \widehat\bigoplus_{a:\,k_a>0}H_a\otimes\Delta\C.
 \tag{MQ16}\label{mq:eq:diagonal-closure}
\]
The attained Hilbert quotient is therefore
\[
 \mathscr M_W=\overline D_W/\overline{\delta W}
 \simeq\widehat\bigoplus_{a:\,k_a\ge2}
          H_a\otimes(\Delta\C)^{\perp}\subset
          \widehat\bigoplus_a H_a\otimes\C^{J_a},
 \tag{MQ17}\label{mq:eq:positive-quotient}
\]
where the displayed last containment is through zero-sum marked
vectors. The natural map
$D_W/\delta W\to\mathscr M_W$ is injective and has dense image.
For $d_a=(d_{a,p})_{p\in J_a}$ put
$\bar d_a=k_a^{-1}\sum_{p\in J_a}d_{a,p}$. Its exact norm is
\[
 \|[d]\|^2=
 \sum_{a:\,k_a>0}\sum_{p\in J_a}\|d_{a,p}-\bar d_a\|^2
 =\sum_{a:\,k_a>0}\frac1{k_a}
           \sum_{\{p,q\}\subset J_a}\|d_{a,p}-d_{a,q}\|^2.
 \tag{MQ18}\label{mq:eq:variance}
\]
The unordered-pair sum counts every distinct pair once.
Thus (CCF31)--(CCF32) give the claimed injective positive
realizations, including every rational-label mixed component.
\end{theorem}
\begin{proof}
On an algebraic channel $W_a$, the map is exactly its diagonal
into its $k_a$ marks. Each $W_a$ is dense in $H_a$, and finite
collections of their diagonal vectors are images of one finite sum
in $W$. Their closure is therefore the indicated Hilbert sum of
diagonals. Conversely that sum is closed: it is the range of the
bounded orthogonal projection taking the arithmetic mean in every
finite marked factor. This proves equality in
\eqref{mq:eq:diagonal-closure}, including its closedness.

The complement of a diagonal vector space consists of vectors with
zero sum. Subtracting $\bar d_a$ gives the corresponding orthogonal
projection, hence proves \eqref{mq:eq:positive-quotient} and the first
norm formula. Expand the squared differences in the second formula:
each squared norm occurs $k_a-1$ times and each distinct inner
product occurs with coefficient $-2\operatorname{Re}$. Dividing
by $k_a$ gives
$\sum_p\|d_{a,p}\|^2-k_a\|\bar d_a\|^2$, which is exactly
the first expression.

An algebraic $d$ has finitely many channels and each coordinate
lies in its corresponding algebraic $W_a$. If its projection is
zero, all its marked coordinates in a channel equal one common
$w_a\in W_a$. The finite sum $w=\sum_a w_a$ then satisfies
$d=\delta w$. Thus
$D_W\cap\overline{\delta W}=\delta W$, proving injectivity.
Finite algebraic zero-sum vectors are dense in the displayed
orthogonal Hilbert sum, proving density. This proof also shows
why changing from evaluation complements to orthogonal complements
does not change injectivity: both decompose the same original
algebraic $D_W$ and the same original diagonal $\delta W$.
\end{proof}

For an explicit nonzero class, choose $p\ne q$ and a nonzero
$h\in V_{\{p,q\}}^\perp$. Give it the $p$ mark and give the
$q$ mark zero. Its projected marked vector is $(h/2,-h/2)$
and its squared quotient norm is $\|h\|^2/2$. There is also a
mixed class involving no nonstandard finite coefficient at all:
take the original label $r=pq$, the channel $T=\varnothing$,
and $h\in\mathcal S(\mathbb R)\otimes\bigotimes_\ell\xi_\ell$.
Its marked set is again exactly $\{p,q\}$ and the same norm
calculation applies to $hU_{pq}$. This verifies the rational-label
part independently of coefficient ramification.

\begin{proposition}[All simultaneous compatibility equations]
\label{mq:prop:compatibility}
For $d=(d_{p,r})\in\overline D_B$, regard $d_{p,r}$ as a vector
of $Q_{p,r}H$. Then
\[
 \overline{\delta B}
 =\bigcap_{r\in\Q^\times}\bigcap_{p\ne q}
   \ker\bigl(d\mapsto Q_{q,r}d_{p,r}-Q_{p,r}d_{q,r}\bigr).
 \tag{MQ19}\label{mq:eq:compatibility}
\]
The unlabelled assertion is obtained by omitting $r$ and using
$Q_p$. Every displayed map is bounded with norm at most $\sqrt2$.
\end{proposition}
\begin{proof}
Coordinate extraction and each $Q$ are contractions. The triangle
and Cauchy--Schwarz inequalities give the stated bound. Commutation
of the $Q$ projections proves that these maps kill every generic
diagonal and hence its closure. Conversely, decompose into the
orthogonal channels. In channel $a=(r,T)$, its $p$ coordinate is
zero outside $J_a$. For $p,q\in J_a$ the displayed equation reads
$d_{a,p}=d_{a,q}$, because both projections act as the identity.
Thus every channel is diagonal. Theorem~\ref{mq:thm:closure} gives
exactly \eqref{mq:eq:compatibility}.
\end{proof}

\subsection{The generic diagonal as an actual closed unbounded operator}

The preceding positivity does not require the generic diagonal to
be a bounded operator in the original generic $L^2$ norm. In fact
it is unbounded. The precise receiving operator is as follows.

\begin{theorem}[Domain, adjoint, and closed range]
\label{mq:thm:closed-diagonal}
Let $\mathscr H_W=\widehat\bigoplus_aH_a$ denote $H$ or $H_B$.
The closure of the algebraic operator $\delta:W\to\overline D_W$
is
\[
 \Dom\delta_{\max}=
 \left\{h=(h_a):\sum_a k_a\|h_a\|^2<\infty\right\},
 \qquad(\delta_{\max}h)_{a,p}=h_a\quad(p\in J_a).
 \tag{MQ20}\label{mq:eq:closed-domain}
\]
Here $h$ is also required to belong to $\mathscr H_W$, so the
unmarked $k_a=0$ part has its usual square summability. Its range
is the closed subspace \eqref{mq:eq:diagonal-closure}. Its adjoint is
\[
 (\delta_{\max}^*d)_a=\sum_{p\in J_a}d_{a,p},\qquad
 \Dom\delta_{\max}^*=
 \left\{d:\sum_a\left\|\sum_{p\in J_a}d_{a,p}\right\|^2
                         <\infty\right\}.
 \tag{MQ21}\label{mq:eq:closed-adjoint}
\]
On an unmarked channel the adjoint value is zero. The positive
operator $\delta_{\max}^*\delta_{\max}$ is the multiplication
operator $N:h_a\mapsto k_ah_a$, with domain
$\sum_a k_a^2\|h_a\|^2<\infty$ in $\mathscr H_W$.
\end{theorem}
\begin{proof}
The squared norm of an algebraic diagonal is exactly
$\sum_a k_a\|h_a\|^2$. Convergence of a vector and its diagonal
implies convergence in each channel; the limiting marked
coordinates are therefore the same coefficient in each mark.
This proves closedness of the operator in
\eqref{mq:eq:closed-domain}. Conversely truncate a vector to finitely
many channels and approximate each remaining $h_a$ by its dense
algebraic tests. The graph norm on that channel is
$(1+k_a)^{1/2}\|h_a\|$, so these approximations converge in
graph norm. Thus this closed operator is exactly the closure of
the original operator, not an unrelated extension.

A vector in \eqref{mq:eq:diagonal-closure} has coefficients $z_a$
repeated $k_a$ times with $\sum k_a\|z_a\|^2<\infty$. Setting
the unmarked coefficients to zero makes $(z_a)$ a vector of
$\mathscr H_W$, since every nonzero $k_a$ is at least one.
It belongs to the displayed domain and maps to that vector.
This proves the closed-range assertion. Pairing against arbitrary
finite-channel tests shows that the adjoint is the marked sum;
this functional is bounded on $\mathscr H_W$ exactly under the
square-summability condition \eqref{mq:eq:closed-adjoint}. Applying
it to a diagonal multiplies by $k_a$. The domains then give
precisely the asserted domain of $N$.
\end{proof}

The restriction to the orthogonal complement of its kernel has a
bounded inverse on the range. There is also a bounded generalized
inverse on all of $\overline D_W$: send a marked vector to
$\bar d_a$ for $k_a>0$ and to zero for $k_a=0$. Indeed
\[
 \sum_a k_a\|\bar d_a\|^2\le\sum_{a,p}\|d_{a,p}\|^2.
 \tag{MQ22}\label{mq:eq:mean-bound}
\]
Thus this mean lies in the actual closed domain and its diagonal
is the orthogonal projection onto \eqref{mq:eq:diagonal-closure}.

For a concrete failure of boundedness, choose finite sets $T_n$
with $|T_n|=n$, and a test $h_n\in V_{T_n}^\perp$ with
$\|h_n\|=1$. Such tests exist: at $p$ the nonzero test
$1_{p\Z_p}-p^{-1}\xi_p$ is orthogonal to $\xi_p$, and its
normalization, together with a normalized real Schwartz function,
gives the required tensor. Then $\|\delta h_n\|=\sqrt n$.
Also $\delta W$ itself need not be closed: choose distinct primes
$p_n$ and unit tests in the corresponding singleton channels.
The convergent marked sum $\sum_n2^{-n}\delta h_n$ lies in
\eqref{mq:eq:diagonal-closure}, but cannot lie in $\delta W$ since
it has infinitely many nonzero channels. These facts verify the
need for the closure used in (CCF30)--(CCF32).

\subsection{The relative completion and the positive arithmetic image}

The algebraic splitting in \eqref{mq:eq:relative-exact} gives the
positive relative receiver chosen in (CCF27): give its $D_W$ the
marked norm and its remaining support groups chosen positive
finite-dimensional forms tensored with the generic $L^2$ norm.
All these forms are strictly positive on the original algebraic
groups. This is a valid specified receiver. It is useful to compute
its exact relation to completing the terms of
\eqref{mq:eq:relative-complex} in their original coefficient norms.

Give support cochains the norm that sums squared coefficient norms
over strict chains. Put $M=|S|>0$. On closed degree-one support
cochains, $c_s$ is constant on each connected component. Write
\[
 \bar c=M^{-1}\sum_{s\in S}c_s,\qquad
 c'_s=c_s-\bar c,\qquad \sum_sc'_s=0.
 \tag{MQ23}\label{mq:eq:support-mean}
\]
The number of points in each support component is thus retained
in the sum; it is not replaced by a component count. The closed
degree-zero boundary is the graph
$w\mapsto(\delta_{\max}w,\Delta w)$ on
$\Dom\delta_{\max}$, and it has closed range because its
second coordinate has norm $\sqrt M\|w\|$.

On an algebraic closed representative define
\[
 q=d-\delta\bar c,\qquad c'=c-\Delta\bar c.
 \tag{MQ24}\label{mq:eq:relative-coordinates}
\]
These are invariant under its actual boundaries. In a channel
with $k_a>0$, write $q_a=q_a^{\circ}+\Delta\bar q_a$, where
$\sum_{p\in J_a}q_{a,p}^{\circ}=0$. Its exact quotient norm
in this closed-graph Hilbert complex is
\[
 \|q_a^{\circ}\|^2+
       \frac{k_aM}{k_a+M}\|\bar q_a\|^2+
       \sum_{s\in S}\|c'_{a,s}\|^2.
 \tag{MQ25}\label{mq:eq:graph-relative-norm}
\]
For $k_a=0$ only the last summand remains. To prove the formula,
all representatives in that channel have the form
$(q_a+\Delta h,\Delta h+c'_a)$. Orthogonality of the two
mean-zero parts leaves the minimization
\[
 \inf_h\left\{\sum_{p\in J_a}\|q_{a,p}+h\|^2
                              +M\|h\|^2\right\}.
\]
Completing the square gives the minimizing value
$h=-(k_a+M)^{-1}\sum_pq_{a,p}$ and the value
$\|q_a\|^2-\|\sum_pq_{a,p}\|^2/(k_a+M)$, which is exactly
\eqref{mq:eq:graph-relative-norm}. Finite algebraic closed
representatives are dense in the full closed degree-one space:
truncate its channels and approximate the coefficient of each
support component by algebraic tests, retaining equality within
that component. Thus taking the orthogonal sum over all channels
proves the formula for the completed receiver. The coordinate $q$
after this completion belongs to the weighted space specified by
\eqref{mq:eq:graph-relative-norm}; its diagonal part need not belong
to the original unweighted $\overline D_W$. This is why no
unjustified application of $\delta_{\max}$ to an arbitrary
completed generic coefficient is required.

The diagonal $q$ coefficient has norm factor $k_aM/(k_a+M)$
in this receiver. On the subspace $c'=0$, the algebraically
split $D_W$ coordinate for every point choice is exactly $q$,
and its norm has factor $k_a$. Since $k_a$ is unbounded, the identity on the
algebraic relative groups does not identify these two completions
by a bounded inverse. This is an exact distinction between two
specified relative norms, not a failure of the arithmetic quotient.
Indeed the relative-to-arithmetic map on either receiver is
\[
 (q,c')\longmapsto q^{\circ}\in\mathscr M_W.
 \tag{MQ26}\label{mq:eq:completed-arithmetic-image}
\]
It is bounded and onto in the closed-graph receiver, by
\eqref{mq:eq:graph-relative-norm}; set $q=q^{\circ}$ and $c'=0$
for an isometric right inverse. On algebraic classes it is exactly
\eqref{mq:eq:arithmetic-map}, because $\delta\bar c$ is an actual
diagonal. Hence both completions retain exactly the same attained
positive mixed-place arithmetic quotient, with all labels and
marked multiplicities. The additional support components and
all higher support cycles remain present in their respective
relative groups.

\subsection{Scaling, Fourier, and the nonunit-prime defect}

The real action $\vartheta_\lambda f(x_\infty,x_f)
=f(\lambda^{-1}x_\infty,x_f)$ commutes with every $P_p$ and
has squared-norm factor $\lambda$. It preserves every orthogonal
channel and its marked set. It therefore gives on the actual
arithmetic quotient
\[
 \langle\vartheta_\lambda[d],\vartheta_\lambda[d']\rangle
       =\lambda\langle[d],[d']\rangle.
 \tag{MQ27}\label{mq:eq:scaling}
\]
This holds for all rational labels as well. It also preserves the
closed diagonal domain, since the weights $k_a$ are unaffected.
The finite-place Fourier transform fixes $\xi_p$ under the
specified conductor-$\Z_p$ characters and Haar measures. Its
unitarity therefore makes it commute with $P_p$: both the line
$\C\xi_p$ and its orthogonal complement are invariant. Full
Fourier is consequently unitary on each marked quotient and
preserves the generic diagonal and its closure. In the crossed
product the map additionally sends $r$ to $r^{-1}$; the equality
$D(r^{-1})=D(r)$ preserves every marked set. These observations
prove that the positive arithmetic receiver in
\eqref{mq:eq:positive-quotient} carries the claimed Fourier and
real-scaling maps. Their exact relation is
$\mathcal F\vartheta_\lambda
=\lambda\vartheta_{\lambda^{-1}}\mathcal F$, obtained by the
real change of variable $x_\infty=\lambda t$.

For comparison, the original rational-prime action
$\alpha_pf(x)=f(p^{-1}x)$ scales all places. The product formula
makes it an isometry of the global $L^2$ space. On an original
unramified coefficient $f=\xi_p\otimes g$, however,
\[
 Q_p\alpha_pf=
 (1_{p\Z_p}-p^{-1}\xi_p)\otimes\alpha_p^{(p)}g,
 \tag{MQ28}\label{mq:eq:prime-defect}
\]
since $P_p1_{p\Z_p}=p^{-1}\xi_p$. Its squared norm is
\[
 (p^{-1}-p^{-2})\|\alpha_p^{(p)}g\|^2
       =(1-p^{-1})\|g\|^2.
 \tag{MQ29}\label{mq:eq:defect-norm}
\]
The other local moduli have product $p$, proving the second
equality. Equations (CCF33) and (CCF34) are thereby verified with their original
normalizations. At $\ell\ne p$ the rational $p$ is a local
unit and preserves the unramified line there. In a crossed-product
stalk $B_p$ all its original labels are units at $p$; conjugation
by $U_p$ leaves those labels unchanged and gives exactly this
coefficient defect. Thus no endomorphism of the $p$-marked
quotient is inferred from a map that fails to preserve its kernel.
The actual defect map and its positive norm remain specified.

\begin{figure}[ht]
\centering
\begin{tikzcd}[column sep=large,row sep=large]
 D_W \arrow[r] \arrow[d,hook]
 & D_W/\delta W \arrow[d,hook] \\
 \overline D_W \arrow[r]
 & \mathscr M_W
\end{tikzcd}
\caption{The exact positive arithmetic receiver. The upper map is
the canonical relative arithmetic image in
Proposition~\ref{mq:prop:arithmetic-image}; the lower map subtracts
the arithmetic mean in every marked channel. Its kernel is the
closed diagonal of Theorem~\ref{mq:thm:closure}. The right-hand map
is injective because the algebraic marked space meets that closed
diagonal exactly in $\delta W$. Formula \eqref{mq:eq:variance}
retains every marked-prime multiplicity and every rational label.}
\end{figure}

The calculations affirm the claimed arithmetic image and attained
positive realization. The additional analytic distinction is the
proved unboundedness of the generic diagonal and the resulting
difference between specified relative completion norms. Its exact
closed domain, range, adjoint, norm comparison, and common arithmetic
quotient have all been constructed above. No purity or zeta-zero
claim is inferred from these positive receiving spaces alone.

